% ======================================================================
% Chapter 2: Standard Model of Particle Interactions
% Source: R. Gupta, "Introduction to Lattice QCD" (arXiv:hep-lat/9807028),
% Section 2, pp. 7-9.
% ======================================================================

\chapter{Standard Model of Particle Interactions}

The Standard Model (SM) of particle interactions is a synthesis of three
of the four forces of nature. These forces are described by gauge
theories, each of which is characterized by a coupling constant as listed
below.

\begin{center}
\begin{tabular}{ll}
STRONG INTERACTIONS & $\alpha_s \sim 1$ \\
ELECTROMAGNETIC INTERACTIONS & $\alpha_{em} \approx 1/137$ \\
WEAK INTERACTIONS & $G_F \approx 10^{-5}\,\mathrm{GeV}^{-2}$.
\end{tabular}
\end{center}

The basic constituents of matter are the six quarks, $u$, $d$, $s$, $c$,
$b$, $t$, each of which comes in 3 colors, and the six leptons $e$,
$\nu_e$, $\mu$, $\nu_\mu$, $\tau$, $\nu_\tau$. The quarks and leptons are
classified into 3 generations of families. The interactions between these
particles is mediated by vector bosons: the 8 gluons mediate strong
interactions, the $W^\pm$ and $Z$ mediate weak interactions, and the
electromagnetic interactions are carried by the photon $\gamma$. The weak
bosons acquire a mass through the Higgs mechanism, and in the minimal SM
formulation there should exist one neutral Higgs particle. This has not yet
been seen experimentally. The mathematical structure that describes the
interactions of these constituents is a local gauge field theory with the
gauge group $SU(3) \times SU(2) \times U(1)$. The parameters that
characterize the SM are listed below; their origins are in an even more
fundamental but as yet undiscovered theory.

\begin{table}[htbp]
\centering
\caption*{Parameters in the Standard Model}
\label{tab:sm-parameters}
\begin{tabular}{l c p{7cm}}
\toprule
Parameters & Number & Comments \\
\midrule
Masses of quarks & 6 & \begin{tabular}[t]{@{}l@{}}
$u$, $d$, $s$ light \\ $c$, $b$ heavy \\ $t = 175 \pm 6$~GeV \end{tabular} \\
Masses of leptons & 6 & \begin{tabular}[t]{@{}l@{}}
$e$, $\mu$, $\tau$ \\ $M_{\nu_e}$, $\nu_\mu$, $\nu_\tau = ?$ \end{tabular} \\
Mass of $W^\pm$ & 1 & 80.3~GeV \\
Mass of $Z$ & 1 & 91.2~GeV \\
Mass of gluons, $\gamma$ & & 0 (Gauge symmetry) \\
Mass of Higgs & 1 & Not yet seen \\
Coupling $\alpha_s$ & 1 & $\approx 1$ for energies $\lesssim 1$~GeV \\
Coupling $\alpha_{em}$ & 1 & 1/137 (=1/128.9 at $M_Z$) \\
Coupling $G_F$ & 1 & $10^{-5}\,\mathrm{GeV}^{-2}$ \\
Weak Mixing Angles & 3 & $\theta_1$, $\theta_2$, $\theta_3$ \\
CP Violating phase & 1 & $\delta$ \\
Strong CP parameter & 1 & $\Theta = ?$ \\
\bottomrule
\end{tabular}
\end{table}

The question marks indicate that these values are tiny and not yet
measured. No mixing angles have been put down for the lepton sector as I
have assumed that the neutrino masses are all zero. A status report on
experimental searches for neutrino masses and their mixing has been
presented by Prof.~Sciulli at this school. The structure of weak
interactions and the origin of the weak mixing angles has been reviewed by
Daniel Treille, while David Kosower covered perturbative QCD.

In the SM, the charged current interactions of the $W$-bosons with the
quarks are given by
\begin{equation}
H_W = \frac{1}{2} \left( J_\mu W^\mu + \text{h.c.} \right)
\label{eq:hw}
\end{equation}
where the current is
\begin{equation}
J_\mu = \frac{g_w}{2\sqrt{2}} \left( \bar{u},\ \bar{c},\ \bar{t} \right)
\gamma_\mu (1 - \gamma_5)\, V
\begin{pmatrix} d \\ s \\ b \end{pmatrix}
\label{eq:jmu}
\end{equation}
and the Fermi coupling constant is related to the $SU(2)$ coupling as
$G_F / \sqrt{2} = g_w^2 / 8 M_W^2$. $V$ is the $3 \times 3$
Cabibbo-Kobayashi-Maskawa (CKM) matrix that has a simple representation in
terms of the flavor transformation matrix
\begin{equation}
V_{CKM} =
\begin{pmatrix}
V_{ud} & V_{us} & V_{ub} \\
V_{cd} & V_{cs} & V_{cb} \\
V_{td} & V_{ts} & V_{tb}
\end{pmatrix}
\label{eq:ckm}
\end{equation}
This matrix has off-diagonal entries because the eigenstates of weak
interactions are not the same as those of strong interactions, i.e.\ the
$d$, $s$, $b$ quarks (or equivalently the $u$, $c$, $t$ quarks) mix under
weak interactions.

For 3 generations the unitarity condition $V^{-1} = V^\dagger$ imposes
constraints. As a result there are only four independent parameters that
can be expressed in terms of 4 angles, $\theta_1$, $\theta_2$, $\theta_3$
and the CP violating phase $\delta$. The representation originally proposed
by Kobayashi and Masakawa is
\begin{equation}
V_{CKM} =
\begin{pmatrix}
c_1 & -s_1 c_3 & -s_1 s_3 \\
s_1 c_2 & c_1 c_2 c_3 - s_2 s_3 e^{i\delta} & c_1 c_2 s_3 + s_2 c_3 e^{i\delta} \\
s_1 s_2 & c_1 s_2 c_3 + c_2 s_3 e^{i\delta} & c_1 s_2 s_3 - c_2 c_3 e^{i\delta}
\end{pmatrix}
\label{eq:km}
\end{equation}
where $c_i = \cos\theta_i$ and $s_i = \sin\theta_i$ for $i = 1, 2, 3$. A
phenomenologically more useful representation of this matrix is the
Wolfenstein parameterization (correct to $O(\lambda^4)$)
\begin{equation}
V_{CKM} =
\begin{pmatrix}
1 - \lambda^2/2 & \lambda & \lambda^3 A (\rho - i\eta) \\
-\lambda & 1 - \lambda^2/2 & \lambda^2 A \\
\lambda^3 A (1 - \rho - i\eta) & -\lambda^2 A & 1
\end{pmatrix}
+ O(\lambda^4)
\label{eq:wolfenstein}
\end{equation}
where $\lambda = \sin\theta_c = 0.22$ is known to 1\% accuracy, $A \approx
0.83$ is known only to 10\% accuracy, and $\rho$ and $\eta$ are poorly
known. The elements $V_{td}$ and $V_{ub}$ of the CKM matrix given in
Eq.~\eqref{eq:wolfenstein} are complex if $\eta \neq 0$, and as a result
there exists a natural mechanism for CP-violation in the SM. The
phenomenology of the CKM matrix and the determinations of the four
parameters have been reviewed by Profs.~Buras and Richman at this school
[10,11].

A list of processes that are amongst the most sensitive probes of the CKM
parameters is shown in Fig.~\ref{fig:1}. Lattice QCD calculations are
beginning to provide amongst the most reliable estimates of the hadronic
matrix elements relevant to these processes. Details of these calculations
and their impact on the determination of $\rho$ and $\eta$ parameters will
be covered by Prof.~Martinelli at this school. My goal is to provide you
with the necessary background.
